\section{IP Architecture}

The Alliance's intellectual-property strategy is the single most
load-bearing element of the compounding thesis. Without enforceable IP,
the substrate is a commodity and the Alliance is a coordination club
without a moat. With enforceable IP --- properly tiered, properly
licensed, and properly defended --- the substrate becomes a perpetual
royalty stream, member admission becomes a controlled gate, and every
Phase-Three acquisition inherits an immediate competitive moat against
its category incumbents.

\subsection{Three-tier licensing}

The substrate ships under the three-tier model summarised in \S 2.1 and
documented in detail at \texttt{LICENSING-POLICY.md} in the Alliance
legal estate:

\textbf{Tier 1 (BSD-3-Clause).} Commodity infrastructure --- node,
codec, database, p2p, SDK, CLI, API. Free for any commercial use,
including by Alliance non-members and by competitors. This tier exists
to maximise the substrate's surface area in the developer ecosystem.

\textbf{Tier 2 (W3A License v1.2, LEL).} Source-available,
patent-protected. Commercial rights granted only to ``Authorized
Networks'' --- the W3A Primary Network, W3A testnets and devnets, and
L1/L2/L3 chains descending from the W3A Primary Network. The
``Descending Chain'' criterion is the operative gate: an Alliance
member chain is automatically authorised; a non-member chain is not.

\textbf{Tier 3 (W3A Research License with Patent Reservation v1.0,
LRL--PR).} Research-only license. Commercial use requires a separate
executed Strategic Commercial License Agreement (SCLA) with W3A
Industries Inc. This tier covers the highest-value primitives where the
Alliance retains direct gate control over commercial use.

A fourth tier --- the closed-source \texttt{\textasciitilde/work/luxcpp/}
private GPU tree --- is licensed only to the Alliance's top settlement
counterparties. It is not source-available in any form.

\subsection{Bifurcated licensor posture (W3A Foundation\ + Hanzo, Inc.)}

The patent estate is bifurcated between two licensors. Lux Industries
Inc.\ licenses the Lux-branded modules (consensus, post-quantum
cryptography, threshold + MPC, FHE, DEX matching, EVM derivatives, VM
substrate, cross-chain messaging). Hanzo, Inc.\ licenses the AI / agent
runtime / MCP / ACI / defense substrate (the AI side of the platform
that the W3A side consumes upstream).

The bifurcation matters for Alliance economics. A Phase-Three
acquisition that takes W3A-substrate IP needs a W3A license. A
Phase-Three acquisition that takes AI / agent IP needs a Hanzo license.
A member that takes both needs both. Two licensors means two streams of
licensing revenue, two layers of defensive posture, and two distinct
SCLA forms --- but always coordinated under the Alliance coordination
layer (\S 8).

\subsection{Member admission as a licensing event}

Every Alliance member's admission is structured as an IP-licensing event:

\begin{itemize}[leftmargin=2em,itemsep=0.1em]
\item Member executes the LEL covering the Tier-2 modules its product
uses.
\item Member executes an SCLA covering any Tier-3 LRL--PR modules.
\item Where the member is taking Hanzo-source modules, member executes
the parallel Hanzo license.
\item Member's chain (or its rollup on the W3A L2) is registered as an
``Authorized Network'' under the LEL Descending Chain criterion.
\item Member's brand integration is registered in the Alliance brand
registry (consistent with the brand-via-env-override discipline).
\item Member's compliance and clean-room obligations are accepted under
the relevant policy documents.
\end{itemize}

Member admission is therefore not just a commercial choice --- it is a
contracted, documented IP-licensing event that materially aligns the
member with the Alliance's compounding economics.

\subsection{Defensive posture}

Where the substrate is used without an executed license, the Alliance
enforces. The licensing terms are deliberately written to provide a
clean cause of action for trade-secret misappropriation, copyright
infringement, patent infringement, breach of fiduciary duty (where
applicable), and unjust enrichment, with statutory and exemplary
damages where the conduct is willful.

This defensive posture is not aggressive litigiousness --- the Alliance's
strong preference is always for partnership-first resolution under the
SCLA framework rather than litigation. But the credibility of the
defensive posture is what makes the partnership offer compelling: a
counterparty who knows litigation is an available alternative will
accept the partnership terms that a counterparty without that knowledge
would walk away from. The defensive posture is the BATNA that makes the
partnership outcome workable.

\subsection{Patent prosecution discipline}

The patent estate is actively managed. Eleven (11) inventions are filed
today; fifty-one (51) are filable and should be perfected to non-
provisional within the standard 12-month window of any invention in
active commercial use; twenty-one (21) defensive publications protect
prior art for novel concepts that are not patent-worthy on their own;
seven (7) trade-secret records carry the highest-value primitives that
are kept out of the patent system entirely.

The discipline is to file early and prosecute aggressively in the US
and key European jurisdictions, with PCT extensions where the
commercial market warrants. Outside counsel review is annual at minimum;
the patent docket is treated as a first-class asset on the Alliance
balance sheet.
